\section{Five candidate paper questions}
\label{sec:questions}
The questions below are candidate contributions, not established novel
results. A final novelty search and citation-chain review should precede
submission. The recommended order is Question 1, an identifiability study
from Question 2, and then the relevant method development.

\subsection{Question 1: Does Dysthe improve prediction beyond fitted NLS?}
\paragraph{Question.}
For which steepness, bandwidth, gauge separation, and noise regimes does
the approved spatial Dysthe equation improve withheld-gauge reconstruction
and downstream prediction compared with both fixed and retuned NLS?

\paragraph{Why this could be a paper.}
This directly develops the Dysthe suggestion in \cite{ehlers2024}.
The scientific issue is whether adding known higher-order physics helps
once a simpler model is given a fair calibration opportunity. A regime
map and a causal explanation of gains or failures are stronger than a
single improved error number.

\paragraph{Hypothesis to test.}
Within the justified envelope regime, Dysthe should be most useful when
the leading corrections are observable but the asymptotic model still
applies. It may offer little benefit at weak nonlinearity, with insufficient
measurements, or outside its validity range. This is an expectation, not
a demonstrated result.

\paragraph{Method.}
First use independently refined \eqref{eq:dysthe} trajectories to validate
implementation. Then use separately labelled HOS-generated observations
or suitable laboratory data to assess model discrepancy. Cross three model
arms---fixed NLS, fitted NLS, fixed Dysthe---with two inference engines:
the same tuned PINN recipe and regularized numerical assimilation around
a differentiable or adjoint spectral solver. Give every arm identical
observations, nuisance-parameter treatment, and tuning opportunities.
Fix envelope extraction and surface reconstruction across arms.

Hold out entire initial wave realizations and entire gauge positions.
Report interpolation between observed gauges separately from propagation
beyond the last observed gauge. A forecast experiment must also specify
which measurement times were available at the forecast origin; using a
complete future time series at an upstream gauge changes the task.

\paragraph{Decisive result and stop condition.}
The target is a reproducible region in which Dysthe lowers unseen
complex-field/surface-wave error or improves calibrated event predictions
at a useful cost. If gains disappear under numerical assimilation, they
may have come from training rather than physics. If the simpler fitted
NLS performs equally well within uncertainty, publish that limit only if
the benchmark and explanation are substantial; do not conceal it by
selecting favorable packets.

\subsection{Question 2: Is the nonlocal correction identifiable from sparse data?}
\paragraph{Question.}
Can the observations distinguish mean-flow effects from carrier calibration,
the unknown inlet field, and local nonlinear terms? Which sensor arrangement
or wave-packet family makes this distinction possible?

\paragraph{Why this could be a paper.}
Predictive success does not imply parameter recovery. A parameter may
absorb model discrepancy or compensate for an uncertain initial field.
The contribution would be an identifiability limit and a practical
observation-design remedy. This is our proposed extension of the inverse
problem, not a published Dysthe identifiability theorem.

\paragraph{Method.}
Before training a neural network, form the observation map
$d\approx\mathcal O(\alpha,z_0)$ using the verified solver, where
$\alpha$ denotes diagnostic coefficients and nuisance calibration
parameters, and $z_0$ denotes an inlet parametrization. Compute scaled
sensitivities
\[
 J_\alpha=\frac{\partial\mathcal O}{\partial\alpha},
 \qquad J_0=\frac{\partial\mathcal O}{\partial z_0}.
\]
Study the singular values of the joint, noise-whitened Jacobian and the
information remaining for $\alpha$ after accounting for $z_0$. A small
singular value signals local confounding; full rank alone does not prove
global identifiability. Supplement it with profile likelihoods,
multi-start estimation, repeated-noise trials, and confidence-interval
coverage for known synthetic parameters.

Compare two gauges with three or more, altered separation, and deliberately
varied packet bandwidth and phase modulation. A sensor-selection criterion
can maximize the smallest relevant singular value or a regularized
information determinant. Compare it with uniform placement and a simple
greedy heuristic. Afterwards test whether the same information gains
improve PINN and numerical assimilation.

\paragraph{Model identity and decisive result.}
The approved physical baseline keeps its coefficient tuple fixed.
Coefficient fitting, removing the nonlocal term, or scaling individual
terms creates separately named diagnostic models, never relabelled
baseline evidence. A useful outcome is an experimentally testable minimum
information condition. If observations cannot distinguish the terms,
report predictive uncertainty rather than a precise inferred coefficient.

\subsection{Question 3: Which intervention prevents long-distance phase failure?}
\paragraph{Question.}
Can propagation-wise training and a correctly evaluated nonlocal residual
maintain phase and focusing-event accuracy at a lower total cost than
conservation penalties or residual adaptivity alone?

\paragraph{Novelty boundary.}
Neither conservation-aware NLS PINNs, spectral differentiation, time
decomposition, nor goal-oriented sampling is new
\cite{noether,cpa,tdd,spectral,goal}. The prospective contribution is a
demonstrated mechanism linking Dysthe's nonlocal coupling, accumulated
phase error, and measurable focusing-event failures.

\paragraph{Method.}
Use structured full-$\tau$ slices so that the mean-flow term is evaluated
by a differentiable Fourier multiplier with controlled padding and
quadrature. Compare local automatic differentiation plus the common
nonlocal operator with a fully spectral profile residual. An optional
weak formulation integrates local derivatives by parts and moves the
self-adjoint multiplier onto a suitable test function where justified.
Check weak residuals against an independent strong residual, since an
incomplete test basis can miss modes.

Start with a compact real-output network and a parameter-matched complex
network. Compare global training, causal weights in $\xi$, and sequential
$\xi$ slabs; distinguish these interventions. Then compare uniform, RAD,
and sensitivity/goal-oriented sampling with equal total budgets, including
candidate scoring and interface work. For a smooth proxy to a maximum,
one possible functional is
\[
 J_\kappa(u)=\kappa^{-1}
  \log\left(\frac{1}{|\Omega|}\int_\Omega e^{\kappa|u|^2}\dd\xi\dd\tau\right),
 \qquad \kappa>0.
\]
Finite $\kappa$ introduces a deliberate approximation; validate the actual
peak height and location afterwards. Do not tune $\kappa$ on the final
test events.

\paragraph{Decisive result and risk.}
Require better unaligned phase/event error versus wall time, not only
lower invariant drift. Independently vary grid, domain, and residual
quadrature. If a benefit vanishes when aliasing or sampler overhead is
counted, it is not a method improvement. A narrow ablation within
Question 1 may be more valuable than a separate architecture paper.

\subsection{Question 4: When does an operator pay for its training?}
\paragraph{Question.}
Can an input-conditioned Dysthe surrogate reliably generalize to unseen
wave packets and achieve a matched-accuracy cost advantage over repeated
numerical solves or inverse fits?

\paragraph{Novelty boundary.}
Generic water-wave PINO reconstruction already exists \cite{wavepino}.
Specify whether the new operator maps initial fields to trajectories
or measurements to reconstructed fields. These are different datasets
and claims. The prospective advance is phase-aware generalization,
controlled physics correction, and calibrated failure detection for the
approved model.

\paragraph{Method.}
Compare a supervised FNO, a physics-informed operator using the exact
residual, and a hybrid that retains an accurate linear propagator and
learns a constrained nonlinear correction. Add a per-instance PINN,
pretrained-and-refined PINN, and numerical solver/assimilation controls.
Split whole initial fields, including all refinements and augmentations.
Reserve distinct bandwidth, steepness, and sensor configurations for
out-of-distribution testing; do not mix them into a single average.

For an operator with total preparation cost $C_{\rm prep}$, inference plus
adaptation cost $C_{\rm op}$, and a classical per-query cost $C_{\rm ref}$,
the simplest break-even estimate is
\[
 Q_\star=\frac{C_{\rm prep}}{C_{\rm ref}-C_{\rm op}},
 \qquad C_{\rm ref}>C_{\rm op}.
\]
Round upward for integer query counts. If the denominator is nonpositive,
there is no speed break-even under this model. Charge data generation,
preprocessing, tuning, training, failed runs, and deployment-specific
refinement consistently. Compare solutions at the same target accuracy.

\paragraph{Decisive result and risk.}
Demonstrate both a useful $Q_\star$ and acceptable held-out phase/event
errors, with calibrated abstention when the input leaves the validated
regime. Ensemble disagreement alone is not calibrated uncertainty:
assess coverage and failure-detection rates on separate trajectories.
A fast model that misses rare focusing events fails the intended task.

\subsection{Question 5: Can a Dysthe prediction have a useful error certificate?}
\paragraph{Question.}
Can one compute a finite-propagation bound on the error of a learned
Dysthe trajectory that is tight enough to remain informative through
nonlinear focusing?

\paragraph{Novelty boundary.}
NLS stability-based bounds and continuous-domain PINN residual
certification already exist \cite{alejo,eiras}. Repeating a generic
Gronwall argument is not a new contribution. The challenge is a computable,
nonvacuous bound for the actual derivative/nonlocal formulation.

\paragraph{Method.}
Begin with the finite-dimensional Fourier discretization
$\dot z=F_N(z)$. Bound the discrepancy of a learned trajectory using its
initial defect and an independently bounded residual. Investigate
interval- or energy-based local stability estimates, possibly per
propagation slab, rather than an unnecessarily large global Lipschitz
constant. Quantify certificate overhead and effectivity
(bound divided by observed error on resolved test cases).

A certificate relative to this ODE is not yet a certificate for the
continuum PDE. A second stage must control Fourier truncation, quadrature,
aliasing, regularity, and the chosen periodic or whole-line setting.
Physical model discrepancy and observation uncertainty are additional
errors, not consequences of the numerical certificate. State explicitly
which of these quantities is bounded rigorously and which is only measured.

\paragraph{Decisive result and risk.}
The target is a bound useful before and after representative focusing
events, with tractable cost. If all rigorous bounds become vacuous long
before those events, a shorter-horizon result may still be valuable, but
an empirical uncertainty band must not be renamed a proof. This is the
highest-risk direction and would benefit from PDE/numerical-analysis
supervision.
